\theory{Elimination theory}{How can unwanted variables be removed from a system of equations without losing the solutions we care about?}{Elimination theory develops algorithms and structural results for projecting polynomial equations onto fewer variables. Substitution, resultants, discriminants, Groebner bases, and cylindrical algebraic decomposition turn a many-variable problem into equations or logical conditions in selected variables.}{Solving polynomial systems, robot kinematics, computer algebra, geometric intersection, invariant theory, quantifier elimination, CAD, and exact symbolic computation.}{Resultants, Groebner bases, and elimination orders remove variables while preserving the polynomial consequences of a system.}

\paragraph{The problem and its history}

Suppose a robot arm has joint angles $x,y$ and its endpoint coordinates are known. The equations may contain both angles, but we want a condition only on the endpoint. Elimination theory removes variables to produce that condition. The new equation is the shadow of the original solution set under projection.

The subject grew from the need to solve polynomial systems. Bezout's theorem bounded the number of solutions in appropriate projective settings. Gaussian elimination solved linear systems by removing variables one at a time. In the nineteenth century, Hermite and Smith normal forms handled linear Diophantine equations and finitely generated abelian groups. Resultants and discriminants supplied polynomial eliminants. \cite{elimination_theory_ref1,elimination_theory_ref2}

Francis Macaulay developed multivariate resultants and U-resultants. Elimination was central in early algebra, then was neglected when Hilbert's non-effective methods became fashionable. Computer algebra revived it: Groebner bases and cylindrical algebraic decomposition made elimination algorithmic again around the 1970s. \cite{elimination_theory_ref3}

\end{historylist}
\section*{3. Theory}
\subsection*{Core theory}
\paragraph{A simple example and the geometric meaning}

Consider
\[
 x+y=3,\qquad xy=2.
\]
The first equation gives $y=3-x$, and substitution yields
\[
 x(3-x)=2,\qquad x^2-3x+2=0.
\]
The two possible values $x=1,2$ then give $y=2,1$. The variable $y$ has been eliminated, but it can be recovered by back-substitution.

For polynomial equations $f_1(x,y)=\cdots=f_r(x,y)=0$, the solution set is an algebraic variety in the $(x,y)$-plane. Projecting it onto the $x$-axis may create a larger set unless the eliminated coordinate can actually be recovered. An elimination polynomial gives a necessary condition, and additional checks are needed to remove extraneous roots introduced by multiplication or closure.

The algebraic object is an ideal $I$ generated by the equations. The elimination ideal
\[
 I\cap k[x_1,\ldots,x_r]
\]
contains all polynomial consequences involving only the variables we keep. The geometric projection of the original variety is related to the zero set of this ideal, with closure issues controlled by algebraic geometry.

\paragraph{Resultants and discriminants}

For two polynomials in one eliminated variable, the resultant is a determinant built from their coefficients. It vanishes exactly when the polynomials have a common root over an algebraic closure, with multiplicities encoded in the determinant. Thus
\[
 \operatorname{Res}_y(f,g)=0
\]
eliminates $y$ from $f(x,y)=0$ and $g(x,y)=0$.

The discriminant of a polynomial is a special resultant between the polynomial and its derivative. It vanishes when two roots collide, so it detects repeated roots and singular parameter values. In geometry it detects when curves become tangent, intersect nontransversely, or acquire a singularity.

Resultants are effective and invariant under many changes of coordinates, but determinant size grows quickly. They can also vanish for reasons that need interpretation, such as a root at infinity in a projective compactification. \cite{elimination_theory_ref3,elimination_theory_ref4}

\paragraph{Groebner bases and elimination orders}

A Groebner basis is a special generating set for a polynomial ideal, adapted to a chosen monomial order. Buchberger's algorithm repeatedly reduces pairs of polynomials until all hidden cancellations are resolved. The computation may be expensive, but the resulting basis makes membership and elimination questions systematic.

The elimination theorem says that if a Groebner basis is computed with an order that places $y$ above $x$, then the basis elements involving only $x$ form a Groebner basis for the elimination ideal $I\cap k[x]$. In practice, a triangular basis can give equations in one variable first, then in two variables, and so on. This resembles solving by substitution but works for nonlinear systems and does not divide by a coefficient that might be zero.

For example, a computer algebra system can eliminate joint variables from robot equations, obtain a polynomial for a possible position, and then solve for the joints. The system may have complex or repeated roots, so physical constraints and numerical verification remain necessary. \cite{elimination_theory_ref1,elimination_theory_ref5}

\paragraph{Logic and quantifier elimination}

Elimination also has a logical meaning. A statement
\[
 \exists y\;[f(x,y)=0\text{ and }g(x,y)>0]
\]
asks whether there is a value of $y$ satisfying conditions for a given $x$. Quantifier elimination replaces the statement by an equivalent condition involving only $x$.

Over algebraically closed fields, polynomial elimination is connected with quantifier elimination for the theory of algebraically closed fields. Over the real numbers, cylindrical algebraic decomposition partitions space into cells on which polynomial signs are constant. It can decide formulas with polynomial equalities and inequalities, though its worst-case complexity is high.

This connection explains why elimination theory matters beyond equations. It turns geometric existence questions into finite logical descriptions and supports exact analysis in computer-aided design, collision detection, and real algebraic geometry. \cite{elimination_theory_ref6}

\paragraph{What the theory depends on and where it is useful}

Elimination theory depends on polynomial rings, ideals, determinants, linear algebra, algebraic geometry, algorithms, and mathematical logic. It helps algebraic geometry by computing projections and closures, invariant theory by producing coordinate-free polynomial relations, and computer algebra by making symbolic solution procedures possible.

In CAD, elimination detects whether two surfaces intersect or become tangent. In computer vision, it removes hidden depth or calibration parameters. In robot kinematics, it produces equations for reachable positions. In solving equations, it turns many variables into a sequence of lower-dimensional problems. Its limits are computational size, extraneous solutions, multiplicities, and the distinction between algebraic closure and the real solutions required by an application.

\section*{4. Important theorems and results}
\begin{enumerate}[leftmargin=*,itemsep=0.35em]

\item \textbf{Bezout's theorem.} Under suitable projective and generic hypotheses, two plane curves of degrees $m$ and $n$ meet in $mn$ points counted with multiplicity. This gives a bound and explains why multiplicity matters in elimination. \cite{elimination_theory_ref2}

\item \textbf{Resultant criterion.} The resultant of two univariate polynomials vanishes exactly when they have a common root over an algebraic closure. \cite{elimination_theory_ref4}

\item \textbf{Elimination theorem.} A Groebner basis in an elimination order contains a basis for the ideal involving only the variables retained. \cite{elimination_theory_ref1}

\item \textbf{Nullstellensatz connection.} Over an algebraically closed field, a polynomial system has no common solution exactly when its ideal contains $1$, so elimination can produce the contradiction $1=0$. \cite{elimination_theory_ref5}

\item \textbf{Quantifier-elimination principle.} In suitable theories, formulas involving quantified variables can be replaced by equivalent formulas without quantifiers; cylindrical algebraic decomposition provides a real-algebraic algorithm for polynomial inequalities. \cite{elimination_theory_ref6}
\end{enumerate}

\theoryapplications
\theoryconnections{elimination theory}{%
\item Linear algebra supplies elimination by row reduction, commutative algebra supplies ideals and Groebner bases, and algebraic geometry interprets elimination as projection of a solution set.
\item Resultants and discriminants provide compact polynomial tests for common roots and singular behavior.
}{%
\item Elimination theory helps solve polynomial systems, robotics constraints, computer-aided design, and algebraic statistics.
\item It turns hidden variables into explicit conditions, but the degree and coefficient growth can make exact computation expensive.
}
\section*{Glossary}
\begin{description}[leftmargin=*,style=nextline,itemsep=0.3em]
\item[\textbf{Resultant.}] A polynomial in the coefficients of two univariate polynomials that vanishes exactly when the polynomials share a common root; it eliminates one variable from a system.
\item[\textbf{Discriminant.}] A special resultant between a polynomial and its derivative that vanishes when two roots coincide, detecting repeated roots and singular parameter values.
\item[\textbf{Gröbner basis.}] A specially chosen generating set for a polynomial ideal adapted to a monomial order, which makes membership and elimination questions systematic and algorithmic.
\item[\textbf{Elimination ideal.}] The subset of an ideal that involves only selected variables, obtained by intersecting the ideal with a smaller polynomial ring; it captures all consequences of the system in those variables.
\item[\textbf{Quantifier elimination.}] A logical procedure that replaces a formula with quantified variables by an equivalent formula without quantifiers, turning existence questions into explicit conditions.
\item[\textbf{Cylindrical algebraic decomposition.}] An algorithm that partitions real space into cells on which the signs of a family of polynomials are constant, enabling decision procedures for formulas involving polynomial equalities and inequalities.
\item[\textbf{Monomial order.}] A rule for ranking polynomial terms that determines the leading term and drives the reduction steps in Gröbner-basis computations.
\item[\textbf{Buchberger's algorithm.}] The algorithm for computing a Gröbner basis by repeatedly reducing pairs of polynomials until all S-polynomials reduce to zero.
\item[\textbf{Nullstellensatz.}] Hilbert's theorem that a polynomial system has no common solution over an algebraically closed field iff its ideal contains $1$.
\end{description}

\section*{7. Sample Questions}
\emph{Exercise reference:} \cite{elimination_theory_ref1}. The cited edition is the source for the exercise set; consult its Exercises section for the official statement and numbering.
\begin{enumerate}[leftmargin=*,itemsep=0.9em]

\item \textbf{Exercise 1.} Compute the resultant of $f(x)=x^2-2$ and $g(x)=x^2-3$.

\emph{Solution.} The Sylvester matrix is $\begin{bmatrix}1&0&-2&0\\0&1&0&-2\\1&0&-3&0\\0&1&0&-3\end{bmatrix}$. Expanding: $\det=1\cdot\det\begin{bmatrix}1&0&-2\\0&-3&0\\1&0&-3\end{bmatrix}-(-2)\cdot\det\begin{bmatrix}0&1&-2\\1&0&0\\0&1&0\end{bmatrix}$. First $3\times3$ det: $1(9-0)-0+(-2)(0+3)=9-6=3$. Second $3\times3$ det: $0-1(0-0)+(-2)(1-0)=-2$. Total: $3-4=-1$. Since $\operatorname{Res}(f,g)=-1\neq0$, the polynomials share no common root, consistent with roots $\pm\sqrt{2}$ and $\pm\sqrt{3}$.

\item \textbf{Exercise 2.} Eliminate $y$ from $x+y=3$ and $x^2+y^2=5$.

\emph{Solution.} From the first equation, $y=3-x$. Substituting: $x^2+(3-x)^2=5$, i.e., $2x^2-6x+4=0$, so $(x-1)(x-2)=0$, giving $x\in\{1,2\}$. Back-substitution yields $y=2$ and $y=1$. The elimination polynomial is $2x^2-6x+4=0$ and the solution set is $\{(1,2),(2,1)\}$.

\item \textbf{Exercise 3.} Compute the discriminant of $p(x)=x^3-x$.

\emph{Solution.} For $x^3+0\cdot x^2-x+0$: $a=1$, $b=0$, $c=-1$, $d=0$. The discriminant is $\Delta=18abcd-4b^3d+b^2c^2-4ac^3-27a^2d^2=0-0+0-4(1)(-1)-0=4$. Since $\Delta=4\neq0$, all three roots are distinct (they are $0,\pm1$).

\item \textbf{Exercise 4.} Find a Gr\"obner basis for $I=\langle x^2+y-1,\;x-y\rangle$ under lexicographic order $x>y$.

\emph{Solution.} From $x-y$, we get $x=y$. Substituting into $x^2+y-1$ gives $y^2+y-1$. The Gr\"obner basis is $\{x-y,\;y^2+y-1\}$. The elimination ideal is $I\cap k[y]=\langle y^2+y-1\rangle$, with roots $y=(-1\pm\sqrt{5})/2$.

\item \textbf{Exercise 5.} Use the Nullstellensatz to determine whether the system $x^2+1=0$, $x^3+1=0$ has a common solution over $\mathbb C$.

\emph{Solution.} The ideals are $\langle x^2+1\rangle$ and $\langle x^3+1\rangle$. Compute $I=\langle x^2+1,x^3+1\rangle$: dividing $x^3+1$ by $x^2+1$ gives remainder $-x+1$, so $I$ contains $x-1$. Then $x-1$ and $x^2+1$ are coprime (the gcd is $1$), so $I=\langle1\rangle$. By the Nullstellensatz, $V(I)=\varnothing$---the system has no common solution over $\mathbb C$.

\end{enumerate}
\section*{8. References}


\begin{thebibliography}{9}
\bibitem{elimination_theory_ref1} D. Cox, J. Little, and D. O'Shea, \emph{Ideals, Varieties, and Algorithms}, Springer, 4th ed., 2015.
\bibitem{elimination_theory_ref2} I. Shafarevich, \emph{Basic Algebraic Geometry}, Springer, 2nd ed., 1994.
\bibitem{elimination_theory_ref3} F. S. Macaulay, \emph{The Algebraic Theory of Modular Systems}, Cambridge University Press, 1916.
\bibitem{elimination_theory_ref4} I. M. Gel'fand, M. Kapranov, and A. Zelevinsky, \emph{Discriminants, Resultants, and Multidimensional Determinants}, Birkhäuser, 1994.
\bibitem{elimination_theory_ref5} D. Eisenbud, \emph{Commutative Algebra with a View Toward Algebraic Geometry}, Springer, 1995.
\bibitem{elimination_theory_ref6} G. E. Collins, "Quantifier elimination for real closed fields by cylindrical algebraic decomposition," in \emph{Automata Theory and Formal Languages}, 1975.
\end{thebibliography}
